% ============================================================
% CENTRALE INHOUDSLIJST — MODULE ANALYSE
%
% #1 = theorie-activity
% #2 = bijhorende oefeningen-activity
% ============================================================


% ============================================================
% PART 1
% ============================================================


\FVDindexpart{Logica — De taal van correct redeneren}

\FVDchapterpair
    {Thema1-Logica/1-Wiskundige-uitspraken}
    {Thema1-Logica/1-Wiskundige-uitspraken-oefeningen}

\FVDchapterpair
    {Thema1-Logica/2-Uitspraken-verbinden}
    {Thema1-Logica/2-Uitspraken-verbinden-oefeningen}

\FVDchapterpair
    {Thema1-Logica/3-Implicatie-equivalentie}
    {Thema1-Logica/3-Implicatie-equivalentie-oefeningen}

\FVDchapterpair
    {Thema1-Logica/Logische-wetten}
    {Thema1-Logica/Logische-wetten-oefeningen}

\FVDchapterpair
    {Thema1-Logica/Uitspraken-negeren}
    {Thema1-Logica/Uitspraken-negeren-oefeningen}

\FVDchapterpair
    {Thema1-Logica/Kwantoren}
    {Thema1-Logica/Kwantoren-oefeningen}

\FVDchapterpair
    {Thema1-Logica/Kwantoren-negeren-en-verwisselen}
    {Thema1-Logica/Kwantoren-negeren-en-verwisselen-oefeningen}

\FVDchapterpair
    {Thema1-Logica/Redeneringen-analyseren}
    {Thema1-Logica/Redeneringen-analyseren-oefeningen}

\FVDchapterpair
    {Thema1-Logica/Van-logica-naar-bewijs}
    {Thema1-Logica/Van-logica-naar-bewijs-oefeningen}

\FVDindexpart{Deelbaarheid en getaltheorie — Structuur in de gehele getallen}


\FVDindexpart{Algebraïsche structuren — Van concrete bewerkingen naar abstracte structuren}